% =============================================================================
\section{L'application KYST}
% =============================================================================
% KYST (Keep Your Surgeries Timelies) : l'outil web qui met le flux tiré entre
% les mains des équipes. Sources : interface/docs (Features.md, Frontend.md),
% interface/ansible/roles/kyst, docs/todo.md.

% Ouverture : l'application est en ligne, le jury peut l'ouvrir pendant la soutenance.
\begin{frame}{KYST est en ligne}
  \begin{columns}[c,onlytextwidth]
    \begin{column}{0.56\textwidth}
      {\Large\bfseries\color{mcmPrimary} sandbox.rezoleo.fr}\\[0.3cm]
      \small
      KYST (\emph{Keep Your Surgeries Timelies}) est l'application web du
      projet. Elle est \textbf{déployée et utilisable dès maintenant}, sur
      ordinateur comme sur téléphone.
      \vspace{0.15cm}
      \begin{exampleblock}{Ce que l'on peut y faire}
        \begin{itemize}
          \item saisir une demande d'intervention en consultation ;
          \item obtenir deux dates, \textbf{justifiées}, compatibles avec le
            bloc et les lits ;
          \item confirmer une date : la vacation et le lit sont réservés.
        \end{itemize}
      \end{exampleblock}
    \end{column}
    \begin{column}{0.40\textwidth}
      \centering
      \qrcode[height=4.2cm]{https://sandbox.rezoleo.fr}\\[0.2cm]
      {\footnotesize Données de démonstration\\synthétiques, aucune donnée réelle}
    \end{column}
  \end{columns}
\end{frame}

% Architecture : trois services, le navigateur ne parle qu'au front.
\begin{frame}{Architecture}
  \small
  \centering
  \begin{tikzpicture}[
      card/.style={draw=McmTerracotta,rounded corners=2pt,fill=white,
                   minimum height=1.25cm,text width=2.0cm,align=center,
                   inner sep=3pt,font=\footnotesize},
      ext/.style={card,draw=McmMuted,dashed},
      flow/.style={{Latex[length=2mm]}-{Latex[length=2mm]},thick,draw=McmTerracotta},
      lab/.style={font=\scriptsize,text=McmMuted,inner sep=1pt}]
    \node[ext]  (user)  at (0,0)    {\textbf{Navigateur}\\ordinateur\\ou téléphone};
    \node[card] (front) at (3.5,0)  {\textbf{Front}\\SvelteKit};
    \node[card] (api)   at (7.0,0)  {\textbf{API}\\FastAPI};
    \node[card] (db)    at (10.5,0) {\textbf{Base}\\PostgreSQL};
    \node[ext]  (ai)    at (7.0,-1.85) {\textbf{Modules IA}\\prédictions,\\ordonnanceur};
    \draw[flow] (user)  -- node[lab,above=3pt] {HTTPS} (front);
    \draw[flow] (front) -- node[lab,above=3pt] {interne} (api);
    \draw[flow] (api)   -- node[lab,above=3pt] {SQL} (db);
    \draw[flow] (api)   -- node[lab,right=3pt] {contrats Python} (ai);
    \begin{scope}[on background layer]
      \node[draw=McmSage,dashed,rounded corners=3pt,inner sep=5pt,
            fit=(front)(api)(db)(ai)] (srv) {};
    \end{scope}
    \node[font=\scriptsize,text=McmSage,anchor=south east] at ([shift={(0,1pt)}]srv.north east)
      {serveur : Podman rootless, Ansible};
  \end{tikzpicture}
  \vspace{0.05cm}
  \begin{columns}[T,onlytextwidth]
    \begin{column}{0.49\textwidth}
      \begin{block}{Séparation des rôles}
        \footnotesize Le navigateur n'appelle \textbf{jamais l'API} : seul le
        serveur du front lui parle. L'API reste l'unique autorité sur les
        droits et les données.
      \end{block}
    \end{column}
    \begin{column}{0.49\textwidth}
      \begin{exampleblock}{Point de branchement de l'IA}
        \footnotesize Les modèles et le solveur se branchent derrière des
        \textbf{interfaces fixes}, sans modifier l'application.
      \end{exampleblock}
    \end{column}
  \end{columns}
\end{frame}

% Interface : le parcours principal, de la consultation à la sortie.
\begin{frame}{Interface : de la consultation à la sortie}
  \small
  \centering
  \begin{tikzpicture}[
      stage/.style={draw=McmTerracotta,rounded corners=2pt,fill=white,
                   minimum height=1.25cm,text width=2.05cm,align=center,
                   inner sep=3pt,font=\scriptsize},
      human/.style={stage,fill=McmSand},
      flow/.style={-{Latex[length=2mm]},thick,draw=McmTerracotta}]
    \node[human] (s1) at (0,0)     {\textbf{Demande}\\saisie en\\consultation};
    \node[stage]  (s2) at (2.75,0)  {\textbf{Prédiction}\\temps en salle,\\séjour, hébergement};
    \node[stage]  (s3) at (5.50,0)  {\textbf{Dates A et B}\\chacune avec\\ses raisons};
    \node[human] (s4) at (8.25,0)  {\textbf{Confirmation}\\par le\\chirurgien};
    \node[stage]  (s5) at (11.00,0) {\textbf{Réservation}\\vacation\\et lit};
    \foreach \a/\b in {s1/s2,s2/s3,s3/s4,s4/s5} \draw[flow] (\a) -- (\b);
    \draw[flow] (s4.south) -- ++(0,-0.45) -| node[pos=0.25,below,font=\scriptsize,text=McmMuted]
      {écarte : nouvelles dates} (s3.south);
  \end{tikzpicture}
  \vspace{0.15cm}
  \begin{columns}[T,onlytextwidth]
    \begin{column}{0.49\textwidth}
      \begin{block}{Pages}
        \footnotesize
        \textbf{Tableau de bord} : lits du jour, remplissage du bloc,
        demandes en attente.\\
        \textbf{Bloc} : semaine par salle, jauge de chaque vacation.\\
        \textbf{Lits} : occupation prévue sur 7, 14 ou 28 jours.\\
        \textbf{Ressources} : salles, unités, chirurgiens.
      \end{block}
    \end{column}
    \begin{column}{0.49\textwidth}
      \begin{exampleblock}{Choix d'interface}
        \footnotesize
        \begin{itemize}
          \item Le chirurgien \textbf{garde la décision} : rien n'est
            programmé sans confirmation.
          \item Jours tendus signalés par un symbole, pas seulement la
            couleur.
          \item Utilisable sur téléphone, en consultation.
        \end{itemize}
      \end{exampleblock}
    \end{column}
  \end{columns}
  \source{En beige : étapes humaines. Le chirurgien peut aussi relancer le
    calcul ou annuler ; la place est alors libérée.}
\end{frame}

% Fonctionnalités : état réel par rapport au cahier des charges.
\begin{frame}{Fonctionnalités et cahier des charges}
  \footnotesize
  \centering
  \renewcommand{\arraystretch}{1.18}
  \setlength{\tabcolsep}{7pt}
  \begin{tabular}{p{6.6cm}p{5.6cm}l}
    \toprule
    Attendu & Dans KYST & État \\
    \midrule
    Donner une date au patient en consultation & Saisie et dates A/B, sur
      mobile & {\color{McmSage}\textbf{Disponible}} \\
    Prédire séjour, temps opératoire, hébergement & Prédictions avec
      intervalles & {\color{McmAmber}\textbf{Démo}} \\
    Proposer une ou deux dates selon bloc et lits & Vérifiées et justifiées
      & {\color{McmAmber}\textbf{Démo}} \\
    Marge pour les urgences & 10\,\% de chaque vacation réservés &
      {\color{McmSage}\textbf{Disponible}} \\
    Suivi du remplissage des lits & Occupation prévue, jours tendus &
      {\color{McmSage}\textbf{Disponible}} \\
    Expliquer, laisser la main au chirurgien & Raisons affichées,
      confirmation & {\color{McmSage}\textbf{Disponible}} \\
    Apprendre de la réalité observée & Durées réelles, précision mesurée &
      {\color{McmSage}\textbf{Disponible}} \\
    Vacations hebdomadaires par spécialité & Ouverture une à une &
      {\color{McmTerracotta}\textbf{Partiel}} \\
    Lissage week-ends et vacances & --- & {\color{McmMuted}\textbf{Prévu}} \\
    \bottomrule
  \end{tabular}
  \source{Démo : le parcours fonctionne de bout en bout, avec des durées de
    démonstration et un ordonnanceur glouton à la place des modèles et du
    solveur.}
\end{frame}

% Sécurité : application, données, infrastructure.
\begin{frame}{Sécurité et confidentialité}
  \footnotesize
  \begin{columns}[T,onlytextwidth]
    \begin{column}{0.32\textwidth}
      \begin{block}{Comptes et sessions}
        \begin{itemize}
          \item Compte validé par un administrateur, \textbf{rôles} vérifiés
            par l'API.
          \item Mots de passe hachés (\textbf{Argon2}), tentatives limitées.
          \item Jetons en cookies \textbf{HttpOnly}, \textbf{Secure},
            invisibles du JavaScript.
          \item Formulaires protégés contre les requêtes intersites (CSRF).
        \end{itemize}
      \end{block}
    \end{column}
    \begin{column}{0.32\textwidth}
      \begin{exampleblock}{Données patient}
        \begin{itemize}
          \item \textbf{Pseudonymisées} : ni nom ni date de naissance, seulement
            l'année et le sexe.
          \item Tableaux de bord \textbf{agrégés}, sans ligne patient.
          \item Chaque décision garde son auteur et sa date.
          \item Aucune donnée réelle dans le dépôt ni sur le serveur.
        \end{itemize}
      \end{exampleblock}
    \end{column}
    \begin{column}{0.32\textwidth}
      \begin{alertblock}{Infrastructure}
        \begin{itemize}
          \item API et base \textbf{non exposées} : pare-feu fermé par défaut.
          \item Seul le reverse proxy TLS joint le front.
          \item Conteneurs \textbf{sans privilèges} (Podman rootless).
          \item Secrets générés au déploiement, hors du dépôt.
        \end{itemize}
      \end{alertblock}
    \end{column}
  \end{columns}
  \source{Prévu : journal d'audit des accès aux données patient, droits limités
    aux demandes de chaque chirurgien.}
\end{frame}

% Reste a faire : d'abord brancher les travaux des autres parties.
\begin{frame}{Ce qui reste à faire}
  \footnotesize
  \begin{columns}[T,onlytextwidth]
    \begin{column}{0.49\textwidth}
      \begin{alertblock}{Brancher l'IA du projet}
        \begin{itemize}
          \item Modèles de \textbf{durée de séjour} et de \textbf{temps en
            salle}, classifieur ambulatoire / conventionnel.
          \item Remplacer l'ordonnanceur glouton par le \textbf{solveur}
            (tabou, recuit, génétique), hors du fil des requêtes.
          \item Relier le coordinateur de \texttt{hospital\_sim} : urgences,
            annulations, salle indisponible.
        \end{itemize}
      \end{alertblock}
    \end{column}
    \begin{column}{0.49\textwidth}
      \begin{block}{Compléter l'outil}
        \begin{itemize}
          \item \textbf{Planning du jour} : ordre de passage et remise en état.
          \item Modèles de vacations hebdomadaires tirés de l'historique.
          \item Aide au lissage des week-ends et vacances.
          \item Droits par chirurgien et journal d'audit.
          \item Migrations de schéma, tests de bout en bout.
        \end{itemize}
      \end{block}
    \end{column}
  \end{columns}
  \vspace{0.05cm}
  \begin{exampleblock}{Message clé}
    \centering
    L'application est prête à \textbf{recevoir les modèles et le solveur} :
    ils remplacent les valeurs de démonstration sans changer l'interface.
  \end{exampleblock}
\end{frame}

% Rappel : la démonstration est accessible maintenant.
\begin{frame}{Essayez KYST}
  \centering
  \vspace{0.2cm}
  \qrcode[height=5.2cm]{https://sandbox.rezoleo.fr}\\[0.45cm]
  {\LARGE\bfseries\color{mcmPrimary} sandbox.rezoleo.fr}\\[0.25cm]
  \small Demandez un accès sur le site : un administrateur de l'équipe le valide.
\end{frame}
